\documentclass[a4paper]{article}     
\usepackage{amsfonts}
\usepackage{amsmath}
\usepackage{amssymb}
\usepackage{graphicx}
\usepackage{cancel}

\newcommand{\asa}{\Leftrightarrow} %als slechts als
\newcommand{\adR}{\Rightarrow} %als dan Rechts
\newcommand{\adL}{\Leftarrow}
\newcommand{\Lh}{\left(} %Links haakje
\newcommand{\Rh}{\right)}
\newcommand{\Lv}{\left[} %Links vierkanthaakje
\newcommand{\Rv}{\right]}
\newcommand{\La}{\left\{} %Links acolade
\newcommand{\Ra}{\right\}}
\newcommand{\Ras}{\right.} %Rechts acolade stelsel (omdat om stelsels te maken heb je een punt nodig)
\newcommand{\Ls}{\left|} %Links (absolutewaarde)streep
\newcommand{\Rs}{\right|}
\newcommand{\pnR}{\rightarrow} %pijl naar Rechts
\newcommand{\limiet}{\lim\limits_}
\newcommand{\schrap}{\cancel}
\newcommand{\schrapb}{\bcancel} %backslash
\newcommand{\schrapk}{\xcancel} %kruis
\newcommand{\vervang}{\cancelto}
\newcommand{\intgrl}{\displaystyle\int} %integraal teken (bij het berekenen van primitieven)

\begin{document}

\noindent \large Auteur: Yoren Collier \\
\noindent \large Studierichting: Informatica\\
\noindent \large Oefeningenreeks 5A nr. 6.7\\

\medskip

\normalsize

$ BI \intgrl_{-2}^2 2^{3x}$\\

Stel: $t = 3x \adR dt = 3dx \asa dx = \dfrac{dt}{3}$\\

Dan worden de grenzen als volgt: ondergrens: $3 \cdot \Lh -2 \Rh = -6$; bovengrens: $3 \cdot2 = 6$\\

$ \adR \intgrl_{-6}^6 2^t \dfrac{dt}{3} = \Lv \dfrac{2^t}{3\ln \Lh 2 \Rh} \Rv_{-6}^6 $\\

$= \dfrac{2^6}{3 \ln \Lh 2 \Rh} - \dfrac{2^{-6}}{3 \ln \Lh 2 \Rh} = \dfrac{64}{3 \ln \Lh 2 \Rh} - \dfrac{1}{64 \cdot 3 \ln \Lh 2 \Rh} = \dfrac{64^2}{192 \ln \Lh 2 \Rh} - \dfrac{1}{192 \ln \Lh 2 \Rh} = \dfrac{4095}{192 \ln \Lh 2 \Rh} = \dfrac{1365}{64 \ln \Lh 2 \Rh}$\\

\begin{thebibliography}{999}
%\bibitem{a} Referentie
\end{thebibliography}
\end{document} 